%% ============================================================
%%  Response to the reviews of COMNET-D-25-04940R2
%%  (Reviewers #4 and #5), for use with a resubmission.
%%  Red [AUTHOR INPUT] markers must be completed before sending.
%% ============================================================
\documentclass[11pt]{article}
\usepackage[margin=2.3cm]{geometry}
\usepackage[T1]{fontenc}
\usepackage{lmodern}
\usepackage{xcolor}
\usepackage{amsmath}
\usepackage{enumitem}
\usepackage[hidelinks]{hyperref}
\usepackage{microtype}
\usepackage{parskip}

\definecolor{rc}{RGB}{0,70,140}
\newcommand{\RC}[1]{\par\medskip\noindent\textcolor{rc}{\textbf{Comment.}}~\textit{#1}\par}
\newcommand{\AR}{\par\noindent\textbf{Response.}~}
\newcommand{\CH}{\par\noindent\textbf{Changes.}~}
\newcommand{\Dt}{\ensuremath{\Delta t}}
\newcommand{\authorinput}[1]{\textcolor{red}{\textbf{[AUTHOR INPUT: #1]}}}

\begin{document}

\begin{center}
{\Large\bfseries Response to Reviewers}\\[4pt]
Previous submission: Computer Networks, COMNET-D-25-04940R2\\[2pt]
\textit{Deferred Asynchronous Post-Quantum Signature Verification for
Low-Latency TLS~1.3 Handshakes}
\end{center}

We thank Reviewers \#4 and \#5 for their careful reading.  Their comments
led to a substantial rewrite and to an entirely new evaluation.  The most
important change is that \textbf{the evaluation has been redone} with the
setup the reviewers asked for: three-certificate hybrid chains,
X25519MLKEM768 key exchange, client and server in separate network
namespaces on disjoint CPU cores joined by an emulated 1500-byte-MTU link,
1000 handshakes per configuration (over 20{,}000 in total, with 43{,}200
deferred verifications timed individually), and a measured bounded worker
pool.  The new results do not reproduce the latency and throughput gains
reported in the previous version, and we have revised our claims
accordingly.  The main changes are:
\begin{enumerate}[noitemsep]
\item New evaluation on OpenSSL~3.5 (Sections~5 and~6).  For lattice-based
  hybrids, deferral does not measurably reduce latency (median change within
  $\pm$0.2\,ms at RTTs of 0--150\,ms) and does not increase throughput; for
  SLH-DSA-128f it saves 1.6\,ms (13.5\%).  \Dt\ is 0.08\,ms at the median
  with a bounded pool.  The earlier 14--21\,ms savings are withdrawn.
\item The contribution is repositioned as a client-side,
  implementation-level scheduling strategy with specified local state
  semantics, not a protocol mechanism (R5.1).
\item The motivation now separates confidentiality (HNDL, addressed by
  hybrid key exchange) from authentication (PKI migration lead time,
  Mosca's inequality), and explains what DAPV does and does not address
  (R4 detailed 1).
\item New Background (PQ schemes, hybrid constructions, TLS~1.3
  authentication modes) and a broadened Related Work section with 15 new
  references (R4.1, R4.2, R4 detailed 2).
\item The latency model is corrected: extra round trips arise from the TCP
  initial congestion window, not from exceeding the MTU.  The new data
  confirm the model's predictions for all five schemes (0, 0, 0, 1 and 2
  extra round trips).
\item The exponential risk model (old Eq.~6, Fig.~11) and the mixed-baseline
  efficiency metric (old Eq.~7, Table~3) are removed; the leakage bound now
  includes the congestion-window limit and uses measured \Dt\ (R5.3, R5.4).
\item No interpolated, extrapolated or projected values remain (R5.8, R5.9).
  Each result opens with its takeaway (R5.6, R5.7).
\item Implementation, drivers and raw data are released (R4.3).
\end{enumerate}

%% ============================================================
\section*{Reviewer \#4}

\RC{W1. Writing style: introduction and background/related work are too
short; decisions such as using threads rather than instruction-level
parallelism (AVX2/AVX-512) are not justified.}
\AR The Introduction has been rewritten and extended (Sections 1.1--1.5),
and the former one-page ``Background and Related Work'' is now two
sections (Sections~2 and~3).  The threads-versus-SIMD question is answered
in the new Section~4.6: the two are complementary.  SIMD shortens one
verification and is already used (liboqs AVX2 code); deferral moves the
remaining time off the critical path.  We also acknowledge explicitly that
for one AVX2 ML-DSA-44 verification ($\approx$0.1\,ms) deferral can save
at most that amount, and identify the cases where synchronous PQ work is
large enough to matter: multi-certificate chains, multi-algorithm
credentials, SLH-DSA and devices without vector units.  Independent
verifications in these cases can be parallelised across cores, which SIMD
cannot do.
\CH Sections 1, 2, 3 and 4.7.

\RC{W2. The literature review is too narrow.}
\AR Section~3 now covers PQ-TLS measurement studies, size-reduction work
(KEMTLS and its variants, intermediate-certificate suppression, mixed
certificate chains), hybrid-signature theory and IETF composite
signatures, asynchronous cryptographic offload in TLS (QTLS/OpenSSL async),
TLS False Start and 0-RTT as precedents for accepting a bounded weaker
guarantee, and formal analyses of deferred authentication.  Table~6
positions DAPV relative to these by the cost addressed and approach.
In revising the bibliography we also found and corrected two entries that
did not match the cited works (the former Cremers2016 and Cheval2022).
\CH Section 3, Table 6, bibliography.

\RC{W3. The implementation was not made available.}
\AR The provider patch, client, server, link emulator, experiment drivers,
raw per-handshake and per-job measurements, and analysis scripts are
released at \authorinput{URL and DOI}.  Section~5.1 describes the
implementation in detail.
\CH Sections 5.1 and 5.2.

\RC{W4. Methodology: client and server on the same machine compete for CPU;
number of repetitions not justified and too few; why multithread if PQC
verification is fast; how many certificates are used; throughput with one
host is hard to trust.}
\AR We agree with every point and have redone the evaluation.
(i)~\emph{Co-location}: client and server now run in separate network
namespaces with their own TCP/IP stacks, pinned to disjoint cores, with
traffic crossing an emulated link rather than loopback (Section~5.3).  We
also measured the co-located placement explicitly.  In neither placement
does DAPV increase the completed-handshake rate (Table~4); the earlier
throughput gains are withdrawn.  The remaining limitation, that both
namespaces share one virtual machine, is stated in Section~6.4.
(ii)~\emph{Repetitions}: 1000 handshakes per configuration at zero RTT and
200 at each other RTT, interleaved in ten rounds; 43{,}200 individually
timed PQ verifications; throughput runs repeated three times.  Differences
are reported with bootstrap confidence intervals.
(iii)~\emph{Why threads}: Section~4.6 discusses this, and the new data
answer it empirically.  For ML-DSA and FN-DSA (37--59\,$\mu$s per
verification) threads bring no benefit, and a thread per job is slower.
Deferral helps only for expensive verification (SLH-DSA-128f, 0.57\,ms).
(iv)~\emph{Certificates}: all experiments use three-certificate chains
(root, intermediate, leaf; the server sends two), so each handshake
performs three PQ verifications, all of which DAPV defers.  A
single-certificate configuration is kept for comparison.
\CH Sections 5 and 6 (new), Tables 2--4, Figures 2--5.

\RC{W5. netem may not trigger real TCP behaviour, local traffic bypasses
the NIC queue, a single certificate does not trigger TCP mechanisms, and
it is unclear whether PQC key exchange is used.}
\AR All addressed.  Traffic no longer uses loopback: each namespace has a
TAP interface with a 1500-byte MTU and segmentation and receive offloads
disabled, so TCP sends real 1500-byte frames and runs its own congestion
control.  A userspace link emulator adds the delay (Section~5.3).  With
three-certificate chains, the TCP mechanisms the reviewer anticipated are
now visible.  P-384\,+\,ML-DSA-65 needs one extra round trip and
SLH-DSA-128f two, exactly as predicted from the initial congestion window
(Table~2, Section~4.7).  X25519MLKEM768 is used in all runs, so both
confidentiality and authentication are post-quantum; X25519 is measured for
comparison (+0.19\,ms for the hybrid KEM).
\CH Sections 4.7, 5.3, 5.4, 6.1.

\RC{W6. Marginal contribution; consider hybrids with more than two
algorithms, which could use multiple verification threads.}
\AR We have narrowed the claims to what the new data support.  The paper
now reports, with a fair baseline, when deferral helps (expensive
verification) and when it does not (lattice schemes on vector-capable
CPUs), together with the provisional-state semantics and the risk analysis.
Three-certificate chains already give three independent verifications per
handshake, which DAPV spreads over two workers.  Algorithm~1 generalises
directly to multi-algorithm credentials, which we list as future work
together with devices without vector units.

\RC{D1. The motivation is HNDL, which concerns confidentiality, while the
proposal concerns authentication; PKI migration is not mentioned; the
hybrid-for-interoperability claim is inaccurate; why hybrid, why
concatenation, what alternatives; fragmentation is motivated but not
addressed.}
\AR All accepted.  Section~1.1 separates HNDL (confidentiality, countered
by hybrid key exchange) from the authentication threat (active
impersonation once a CRQC exists) and explains why authentication must
nevertheless migrate now: PKI lead times and Mosca's inequality.
Section~1.2 corrects the statement about interoperability: hybrid
authentication is a hedge against breaks of the young PQ schemes, not an
interoperability mechanism; a composite certificate is not verifiable by
clients that do not support it.  Section~2.2 describes the four hybrid
constructions (composite/concatenation, nested extensions, parallel
chains, mixed chains), their trade-offs, and which of them DAPV applies
to.  The abstract and introduction now state that DAPV addresses only the
computational component of the cost and not size or fragmentation.
\CH Abstract, Sections 1.1, 1.2, 2.2.

\RC{D2. Background does not explain PQC, available algorithms, how to
choose, what hybrids are, or how TLS authentication works (not only
certificates; 0-RTT, PSK; Certificate message also carries signatures).}
\AR New Section~2.1 describes ML-DSA, FN-DSA and SLH-DSA, the assumptions
behind them and the selection trade-offs, with a size table (Table~1).
Section~2.3 describes the three TLS~1.3 authentication modes
(certificate-based, PSK/resumption with 0-RTT, post-handshake), notes
that chain validation also requires signature verifications in the
\texttt{Certificate} message, and explains how DAPV interacts with each
mode.  A new rule (Section~4.3) prevents a provisional session's
tickets from being used for resumption or 0-RTT before promotion.
Section~2.3 now states that the implementation defers the PQ components of
the chain signatures as well as of \texttt{CertificateVerify}.

\RC{D3. Section 3 needs more discussion; some messages of Fig.~1 are not
explained; design goals.}
\AR Fig.~1 has been redrawn and now shows all TLS~1.3 messages; every
numbered step is explained in Section~4.2.  Design goals and non-goals are
listed in Section~4.1.

\RC{D4. The key-exchange algorithm is not stated; ML-KEM adds bytes and
computation.}
\AR X25519MLKEM768 is now used in all experiments (Section~5.4), and
X25519 is measured for comparison.  Table~1 gives the X25519MLKEM768 share
sizes (1216/1120 bytes).  The reviewer's figure of 800 bytes corresponds to
the ML-KEM-512 public key.

\RC{D5. Algorithm 1 dispatches the PQC verification after verifying ECDSA.
Why not dispatch first to run in parallel?  Algorithms are not explained.}
\AR Explained in Section~4.5 (``Dispatch order''): dispatching only after
$\sigma_E$ verifies prevents an attacker without a valid classical
signature from consuming workers, and the classical verification takes
only tens of microseconds, so the difference in \Dt\ is negligible.
Implementations may dispatch first and cancel on failure.  All three
algorithms are now explained in the text, and the functions they use are
defined.

\RC{Minor: ``carries'' should be ``carry''.}
\AR The sentence has been rewritten.

%% ============================================================
\section*{Reviewer \#5}

\RC{1. Is DAPV a protocol-level contribution or an implementation-level
scheduling optimisation?}
\AR It is implementation-level; we agree with the reviewer.  DAPV changes
no message and needs no server support.  All ``protocol-level'' claims have
been removed; the Introduction (Section~1.4) states the positioning
explicitly, and the standardisation discussion (Section~8) now says that no
TLS extension is required, only implementation guidance on the provisional
semantics.

\RC{2. What happens in the transition from provisional to
fully\_authenticated?  Is there re-keying?}
\AR No.  Section~4.3 (``Promotion does not re-key'') explains that the
transition is a local state change, why re-keying is unnecessary (the
keys depend only on the key-exchange secret and the transcript), and what
this means for data sent during \Dt\ in both the success and failure
cases.

\RC{3. Justify Equation (6): where does $\beta$ come from, and what event
does $R(\Delta t, B)$ represent?}
\AR On re-examination we could not justify $\beta$ empirically, and the
exponential form did not correspond to a well-defined event.  We have
removed Eq.~(6) and its figure.  Exposure is now characterised only by
the deterministic bound of Eq.~(3), which we have strengthened to include
the TCP congestion-window limit.

\RC{4. Justify Equation (7), which mixes baselines.}
\AR We agree the mixed baselines were not meaningful.  The metric and its
table have been removed; the results report latency and rate changes
directly against a single, consistent baseline: synchronous verification
of the same configuration (Tables~2 and~4).

\RC{5. Justify the 8\,KB cap.}
\AR It is a heuristic default, not an empirical result.  Section~7.2
(``Choice of $D_\text{policy}$'') explains that the correct value depends on
what data the application will reveal to a classical-only authenticated
peer, and gives the reasoning behind 8\,KB as a default.

\RC{6. Clarify the purpose of Tables 2 and 4.}
\AR Old Table~4 has been removed because most of its entries were
interpolated or extrapolated (see comment~8).  The evaluation is now built
around three tables, each introduced by its takeaway: Table~2 (latency,
showing that deferral helps only for SLH-DSA, and that round trips follow
the congestion-window model), Table~3 (\Dt, showing that the provisional
window is a fraction of a millisecond with a bounded pool), and Table~4
(throughput, showing that DAPV does not raise the rate in either CPU
placement).

\RC{7. Improve the presentation of the evaluation.}
\AR The evaluation section has been rewritten around the new
experiments.  It now has three tables and four figures, each introduced by
a one-sentence takeaway.  All eleven previous evaluation figures and the
old efficiency and interpolated tables have been removed.

\RC{8. Why were interpolation and extrapolation used?}
\AR They should not have been.  All values in the revised paper are
measured, at all five RTTs, for every configuration (Table~2 and
Figure~3).

\RC{9. Was the $\Delta t < 1$\,ms thread-pool result implemented or
projected?}
\AR It was a projection in the previous version.  We have now implemented
and measured a bounded pool.  \Dt\ is 0.079\,ms at the median and 0.42\,ms
at the 99th percentile for ML-DSA-44, versus 0.114 and 0.99\,ms with a
thread per job (Table~3, Figure~4).

\RC{10. VerifyPQC() is not defined; unresolved references in Section 5.7.}
\AR VerifyPQC and VerifyECDSA are defined in Section~4.5, including the
message they verify.  The unresolved references (to an undefined
equation label and figure label) have been fixed; the revised manuscript
compiles without undefined references.

\medskip
We thank both reviewers again.

\end{document}
